\subsection*{Scalar Matmul (No DMA)}

The first version uses \texttt{parallel} to divide the output and \texttt{.at()} to read individual elements. It does not yet use explicit data movement.

\begin{lstlisting}[language=croqtile,caption={Scalar Matmul (No DMA)}]
__co__ s32 [128, 256] matmul(s32 [128, 256] lhs, s32 [256, 256] rhs) {
  s32 [lhs.span(0), rhs.span(1)] output;
  parallel {p, q} by [16, 64] {
    foreach {m, n} in [128 / #p, 256 / #q]
      output.at(p#m, q#n) = 0;
    foreach index = {m, n, k} in [128 / #p, 256 / #q, 256]
      output.at(p#m, q#n) += lhs.at(p#m, k) * rhs.at(k, q#n);
  }
  return output;
}
\end{lstlisting}

\begin{figure}[htbp]\centering
\includegraphics[width=\linewidth,height=0.56\textheight,keepaspectratio]{f05-fig5_scalar_matmul_light.png}
\caption{The [128, 256] output has a [16, 64] tile grid. Each tile has shape [8, 4]; only the first few tiles are drawn.}
\end{figure}

The tile dimensions follow from the parallel bounds:

The output declaration derives 128 rows from \co@lhs.span(0)@ and 256 columns from \co@rhs.span(1)@. The contraction dimension appears in both inputs but does not become an output dimension.

The declaration \co@parallel {p, q} by [16, 64]@ gives one parallel level two coordinate axes. Its 1,024 index pairs partition the output into tiles of eight rows and four columns.

The loop binds \co@m@, \co@n@, and \co@k@ as components of \co@index@. The first two visit positions inside the output tile; the third traverses the contraction dimension. Their trip counts follow directly from the output partition:

\begin{itemize}[leftmargin=1.6em]
\item \texttt{128 /\allowbreak{} \#\allowbreak{}p = 128 /\allowbreak{} 16 = 8} rows per tile
\item \texttt{256 /\allowbreak{} \#\allowbreak{}q = 256 /\allowbreak{} 64 = 4} columns per tile
\item \texttt{256} for the full contraction dimension
\end{itemize}

The composed row \co@p#m@ has value \co@p * 8 + m@. As p visits the 16 row tiles and m visits the eight rows within each tile, it covers rows 0 through 127. The column expression \co@q#n@ works the same way with four columns per tile.

For each output position, the body sums \co@lhs.at(p#m, k) * rhs.at(k, q#n)@ over K. All contributions go to the same output element. The zeroing loop runs in the same parallel instance, so each accumulator element is cleared by the instance that owns it.
